\documentclass[11pt,a4paper]{article}
\usepackage[margin=2.2cm]{geometry}
\usepackage{hcqa-style}

\hypersetup{
  pdftitle={hi\_cube\_qa Quickstart},
  pdfauthor={GAMA G23 MeerKAT HI Survey}
}

\pagestyle{empty}

\begin{document}

% --------------------------- title block ----------------------------------
\begin{center}
  {\hcqalogo[1.4]}\\[10pt]
  {\LARGE\bfseries\ttfamily hi\_cube\_qa}\\[4pt]
  {\large Quickstart}\\[6pt]
  {\normalsize MeerKAT \textsc{Hi}\ cube quality assessment --- GAMA~G23}
\end{center}
\noindent\rule{\textwidth}{0.6pt}
\vspace{4pt}

\noindent
\code{hi\_cube\_qa} reads the spectral-line cubes that \textsc{CARACal}
\citep{jozsa2020} writes for an observation session and produces \textbf{one
interactive HTML QA report per session}. It auto-discovers the cubes, measures
the per-channel noise, compares it to the radiometer-equation expectation,
checks the synthesised beam and the integrated spectrum, and grades each cube
against a fixed expert checklist. For the full scientific justification of every
metric, see the companion \emph{Reference Manual} (\code{hi\_cube\_qa.tex}).

% ---------------------------------------------------------------------------
\section*{1.\quad The three-step workflow}
\addcontentsline{toc}{section}{Workflow}

\textbf{Step 1 --- point a config at one or more session roots.}
The input is a list of \emph{sessions}, not cubes; everything under each session is discovered
automatically.
\begin{lstlisting}
sessions:
    - root: /data/G23/1610268336_sdp_l0/caracal
products: [dirty, image, residual]   # default three; or 'all'; or a subset
\end{lstlisting}

\textbf{Step 2 --- preview what will be assessed} (fast, no heavy compute):
\begin{lstlisting}
hi-cube-qa --config qa_G23_all_sessions.yaml --list
\end{lstlisting}
This prints every \code{pointing/cube\_N} unit and the products found, to
confirm e.g.\ that residuals are picked up before committing to a full run.

\textbf{Step 3 --- run it.} Locally:
\begin{lstlisting}
hi-cube-qa --config qa_G23_all_sessions.yaml
\end{lstlisting}
or on the cluster (Slurm + Singularity):
\begin{lstlisting}
sbatch run_qa.sh                  # uses qa_G23_all_sessions.yaml
sbatch run_qa.sh my_config.yaml   # or a custom config
\end{lstlisting}
Each session yields \code{<obsid>\_qa\_report.html} plus three CSVs and a log in
\code{output.dir}, and a master \code{qa\_index.html} that tabs through all
sessions. \textbf{Open \code{qa\_index.html}} in any modern browser.

% ---------------------------------------------------------------------------
\section*{2.\quad What the report contains}
\addcontentsline{toc}{section}{Report contents}

The report is organised by \emph{pointing}; each \code{(pointing, cube\_N)} unit
overlays its products (which share a WCS and channel grid) on common axes:

\begin{description}[leftmargin=1.4em,style=nextline,itemsep=1pt]
  \item[Provenance strip] Briggs robust, niter, restoring beam read from the
        header; a warning if the cube is primary-beam corrected.
  \item[Summary table] Median RMS, RMS/theory, $3\sigma$ $N_{\mathrm{HI}}$,
        channel grades, beam.
  \item[RMS curves] Per-channel noise coloured by quality, with the per-channel
        theory curve (frequency- and flagging-aware); plus RMS/theory, the
        cumulative (converging) RMS, and the RMS histogram.
  \item[Noise overview] Boresight (best-case), median and peak noise versus the
        flat and flagging-reduced theory curves.
  \item[Spectral profile \& band spectrum] Peak/mean per channel, and the
        integrated flux density $S_{\mathrm{int}}(v)$ with $\pm1/\pm3\sigma$
        envelopes, a velocity/frequency dual axis, and the redshift marker.
  \item[$N_{\mathrm{HI}}$ sensitivity] Per-channel $3\sigma$ column-density limit.
  \item[Non-Gaussianity] Per-channel kurtosis (expectation 3 for Gaussian noise)
        -- catches contamination the RMS check alone cannot see.
  \item[Beam statistics] \textsc{Bmaj}/\textsc{Bmin}/\textsc{Bpa} across the band.
  \item[Spatial maps] Per-pixel noise map and continuum-persistence map
        (continuum-subtraction quality) for the residual.
  \item[Channel maps] Diagnostic channels: cleanest/typical/noisiest by RMS and
        strongest/typical/weakest by SNR.
  \item[Checklist] Pass/warn/fail against eleven expert criteria, including baseline flatness and large-scale noise.
  \item[Flagging] One session-level bar chart of \textsc{CARACal} flag fractions.
\end{description}
A master \code{qa\_index.html} tabs through every session, and
\code{qa\_survey\_summary.\{csv,json\}} rolls up every cube across all sessions.

% ---------------------------------------------------------------------------
\section*{3.\quad Three things to set correctly}
\addcontentsline{toc}{section}{Key settings}

\begin{enumerate}[leftmargin=1.6em,itemsep=2pt]
  \item \textbf{\code{radiometer.t\_total}} --- the \emph{on-source} hours. It
        scales the theoretical-noise line $\sigma_{\mathrm{theory}}\propto
        \tau^{-1/2}$, so set it to the real track length or the RMS/theory ratio
        is meaningless. (Channel width is auto-read from each FITS header.)
  \item \textbf{\code{radiometer.nfac}} --- the imaging-weight penalty
        (e.g.\ $\approx 1.7$ for Briggs robust $=-0.5$), or \code{auto} to read
        the robust value from the header; makes the theory match the weighting
        used. See also \code{sefd\_curve} and \code{flag\_fraction}.
  \item \textbf{\code{products}} --- defaults to \code{[dirty, image,
        residual]}. Use \code{all} to also pick up \code{model}/\code{psf}, or a
        subset to keep reports light. The \code{residual} cube is used as the
        noise reference when present.
\end{enumerate}

\vfill
\noindent\rule{\textwidth}{0.4pt}\\[2pt]
{\footnotesize Built on \code{numpy}, \code{scipy}, \code{astropy},
\code{matplotlib} and \code{plotly}. Full method details and references:
\emph{hi\_cube\_qa Reference Manual}.}

\bibliographystyle{plainnat}
\renewcommand{\refname}{\vspace{-1.5em}}
{\footnotesize\bibliography{refs}}

\end{document}
